\abstract % Структурный элемент: РЕФЕРАТ. Строка со сведениями об объёме формируется автоматически.

\keywords{БОЛЬШИЕ ДАННЫЕ, APACHE PAIMON, КОЛОНОЧНЫЙ ФОРМАТ, VORTEX, PARQUET, ORC, LAKEHOUSE, LSM-ДЕРЕВО, КОДИРОВАНИЕ ДАННЫХ, НАГРУЗОЧНОЕ ТЕСТИРОВАНИЕ, APACHE FLINK}

Объект исследования --- колоночные файловые форматы (Parquet, ORC, Vortex) в
составе lakehouse-системы Apache Paimon; предмет --- алгоритмы кодирования,
сжатия и пропуска данных и их влияние на производительность. Цель --- определить
область применимости формата Vortex в Apache Paimon, разработать модуль его
интеграции и провести сравнительный анализ.

Выполнен обзор lakehouse-архитектур и колоночных форматов; разработан и
интегрирован модуль поддержки Vortex (\texttt{FileFormat}, конвертер предикатов,
JNI); проведено функциональное тестирование в контейнерах; развёрнуто кластерное
окружение средствами Ansible; проведены две серии нагрузочных экспериментов ---
batch-аналитика TPC-DS (\SI{100}{\giga\byte}) и near-real-time-сценарий с
потоковой записью и параллельными запросами.

Результаты: на batch-аналитике Vortex --- быстрейший из четырёх форматов
(быстрее всех на \num{90} запросах из \num{104}); по суммарному времени опережает
Parquet на \SI{25}{\percent}, ORC --- на \SI{32}{\percent}, пиковое ускорение ---
до \SI{83}{\percent} (основная область применимости); Lance самый медленный. На NRT-сценарии Vortex
эквивалентен Parquet и ORC по записи; по медиане чтения --- быстрейший формат в
обеих фазах (9/13 запросов), особенно на запросах с селективным предикатом,
диапазонным фильтром по строке-префиксу и в многошаговой аналитике; на простой
агрегации без предиката незначительно уступает (отсутствие проталкивания
агрегатов в интерфейсе Paimon). ORC показывает наиболее стабильное время отклика и занимает меньше всех места
на диске (на 5--18\,\% меньше Parquet и Vortex). Parquet ловит
эпизодические выбросы времени запросов в моменты сброса по чекпойнту, что
согласуется с архитектурным выигрышем Vortex за счёт независимого нативного
пула соединений с хранилищем. Область применения --- проектирование витрин данных
на Apache Paimon, развитие его экосистемы.
